\section{A stronger obstruction: nonreal zeros arbitrarily near the real axis}
\label{sec:nearreal}
Theorem~\ref{thm:finitedepth} exhibits
nonreal zeros that recede into the complex plane
as $N$ grows. That distance restriction is
not essential. The Gaussian parameter can be
chosen to conceal \emph{arbitrarily near-real}
complex zeros from every prescribed finite
number of global Laguerre tests.

\begin{theorem}[Near-real finite-depth barrier]
\label{thm:nearreal}
For every integer $N\ge1$ and every
$\varepsilon>0$ there exists a positive,
even, rapidly decreasing real kernel $K$
whose Fourier transform $F$ has nonreal
zeros at distance less than $\varepsilon$
from the real axis, although
\[
L_n[F](x)>0
\qquad\text{for every real }x,\quad
1\le n\le N.
\]
More explicitly, choose any $d\in(0,1)$,
set
\begin{equation}
A=\frac{16N}{(1-d)^2},\qquad
K(u)=e^{-u^2/(2A)}
\left[1+d\,e^{-1/(2A)}\cosh(u/A)\right],
\label{eq:nearK}
\end{equation}
and define $F(z)=\int K(u)e^{izu}\,du$.
Then
\[
F(z)=\sqrt{2\pi A}\,e^{-Az^2/2}
(1+d\cos z)
\]
and for all real $x$ and $1\le n\le N$
\begin{equation}
\boxed{
L_n[F](x)\ge
\frac{3\pi A}{2}
\,e^{-Ax^2}\,
\frac{(1-d)^2A^n}{n!}>0.
}
\label{eq:nearbound}
\end{equation}
The nonreal zeros are precisely
\[
(2k+1)\pi\pm i\operatorname{arcosh}(1/d).
\]
As $d\uparrow1$, their imaginary parts
tend to zero.
\end{theorem}

\begin{proof}
The inverse Fourier formula follows by
completing the square in the Gaussian
integral, with the $\cosh(u/A)$ term
corresponding to shifts $z\mapsto z\pm i/A$.
The resulting kernel is positive and
rapidly decaying for real $u$.

Fix $1\le n\le N$, and put
\begin{align*}
a&=\frac{A^n}{n!},\\
b&=\sum_{j=1}^n
\frac{A^{n-j}}{(n-j)!(2j)!},\\
c&=\sum_{j=1}^n
\frac{4^jA^{n-j}}{(n-j)!(2j)!}.
\end{align*}
The coefficient of $y^{2n}$ in
$e^{Ay^2}(1+d\cos(x+iy))
(1+d\cos(x-iy))$ is
\begin{equation}
a(1+d\cos x)^2
+2d\,b\cos x+\frac{d^2}{2}c.
\label{eq:nearcoefficient}
\end{equation}
Here $c\ge0$. The other two terms have
the lower bound
\[
a(1-d)^2-2db
\]
because $-1\le\cos x\le1$.
To estimate $b$, use
\[
\frac ba=
\sum_{j=1}^n
\frac{n!}{(n-j)!(2j)!A^j}
\le\sum_{j=1}^n\left(\frac nA\right)^j.
\]
Since $A\ge16N\ge2n$, the
geometric series is bounded by
\[
\frac ba\le
\frac{n/A}{1-n/A}\le\frac{2N}{A}.
\]
Therefore
\[
2db\le\frac{4dN}{A}a
=\frac{d(1-d)^2}{4}a.
\]
Using $d<1$ gives
\[
a(1-d)^2-2db
\ge\frac34a(1-d)^2>0.
\]
Multiplying by the positive Gaussian
factor $2\pi A\,e^{-Ax^2}$ proves
\eqref{eq:nearbound}.

Finally, take for instance
$d=1-m^{-2}$ with integer $m\ge2$.
Then $\cosh y=1/d$ at the nonreal
zero height $y>0$, and
$1/d=1+1/(m^2-1)$.
Since $\cosh\varepsilon\ge
1+\varepsilon^2/2$, any
$m$ satisfying
$m^2-1>2/\varepsilon^2$
gives $0<y<\varepsilon$.
This explicitly establishes the
arbitrary-nearness statement.
\end{proof}

\paragraph{Certification.}
The standard-library script
\texttt{certificates/near\_real\_finite\_depth\_exact.py}
checks the exact rational coefficient
bounds for several pairs $(N,m)$,
including a height threshold below
$10^{-3}$. The companion Lean module
\texttt{XiKernel/NearRealFiniteDepth.lean}
proves the underlying algebraic
positivity gate under the explicit
coefficient hypotheses; the
all-$N$ binomial comparison and
Gaussian integral are given above
as exact paper proofs rather than
end-to-end Lean theorems.

\paragraph{Research significance.}
Even global validity of any
\emph{finite initial block}
of generalized Laguerre inequalities
cannot detect nonreal zeros arbitrarily
near the real axis in this positive
Fourier-kernel class.
The issue is not the distance of a
putative off-line zero, but uniform
control of all Laguerre orders.
This does not decide any zero of
the genuine Riemann $\Xi$ function.
